\chapter{66}

\section{Donnerstag, 30. August}



\emph{Und vielleicht kannst du dein Leiden abkürzen, indem du dich stellst.}

Noch bevor Louis die Augen aufschlug, fiel ihm Valentinas Ratschlag ein. Mehr Schlag als Rat, dachte er sich und 
schüttelte konsterniert den Kopf.  

Erstmal musste er sich den Schafen zu widmen und die Hütte auf Vordermann 
bringen. Bereits am nächsten Morgen würden Joh und Manuela auf die Alp
hochkommen. Zusammen mit David und einer Besucherin. Sie sollten sich
willkommen fühlen. Joh und David würden sich in Mazis Hütte
einquartieren. Er wollte seine Arbeitgeber von sich überzeugen, hatte
aber auch Respekt davor, bald von ihnen allen umgeben zu sein und die
Zeit mit sich alleine zu verlieren. Unten auf dem Hof stand
anschliessend die Zusammenarbeit mit Mazi an.

\sectionbreak

Louis schwang seinen Rucksack über die Schulter. Er lag in der Zeit. Die
Schafe wollte er heute in gleicher Reihenfolge besuchen wie am Vortag
und vermutete, heute schneller voranzukommen. Die würzige Luft tat gut.
Der letzte und kurze Aufstieg zu den Felsen liess ihn sein Tempo
drosseln. In der Mitte blieb er stehen. Er stützte seine Hände auf den
Hüften ab. Am liebsten hätte er seinen Blick eins zu eins festgehalten, als Bild geknipst und mit allen geteilt,
die ihm lieb waren. Es schmerzte, abgeschnitten und alleine hier zu
stehen. Inmitten dieser vollkommenen Schönheit. Selbst im Kontakt mit
Valentina musste er vorsichtig sein. Eine Aufnahme von hier würde sie ahnen
lassen, wo er sich befand. Mit Valentina hatte er
wenigstens einen Anker gefunden. Jemanden, dem er sich ein Stück weit
anzuvertrauen gewagt hatte.

\sectionbreak

Louis wunderte sich einmal mehr, mit welcher Selbstverständlichkeit die Tiere ihn empfingen, wie flauschig sich ihre Felle
anfühlten und wie richtig er sich in ihrer Nähe vorkam.
Der Moment liess eigentlich keine Zweifel zu. Unterdessen aber wusste er, seine Verlorenheit wartete bereits irgendwo. Und im Handumdrehen war sie da. Mit ihr tat sich eine weit zurückliegende Szene mit Giorgia auf, die Louis nicht hatte erleben wollen.
Schon zum zweiten Mal an diesem Morgen hatte sie «Ilarias»
\emph{«Giudizi universali»}  in der Küche laufen lassen. Überlaut, während sie die Spülmaschine ausräumte. Louis hatte ihr dabei erschrocken zugesehen. Ihre schnellen Bewegungen hatten ihn irritiert. Sein Versuch, sich unbeteiligt zu stellen, misslang. 
Er hatte schliesslich ihre Worte über sich ergehen lassen und sich insgeheim gewünscht, sie möge ganz schnell aufhören. Sie hatte zum wiederholten Mal Kritik an ihrer Zweisamkeit angebracht.

Louis wischte sich über die Augen. Blitzartig fühlte er sich 
zurückversetzt in seine abgrenzende Verweigerung, die Giorgia zunehmend 
zur Weissglut getrieben hatte. Und es war, als höre er ihre Worte
erneut, als werfe sie ihm mit Wucht ihre Überzeugungen entgegen. Er hatte längst gewusst, wieviel ihr das Lied bedeutete.
Von der Suche nach Freiheit und der Unabhängigkeit von der Meinung
der Gesellschaft hatte sie gesprochen. Von Erwartungen und Urteilen und dem freien
Willen. Von der Wichtigkeit, das Leben selbstbestimmt zu gestalten und
von irrer Anpassung. Und davon, dass sie beide zusammen feststeckten. 

«In Beziehungen geht es doch ums Gleiche, auch darum, zu wissen, was
jeder will, in Freiheit verbunden zu sein, Abhängigkeiten zu erkennen,
aufzulösen und gemeinsam weiterzugehen - ich hab den
Eindruck, du passt dich mir einfach an, Louis - ich weiss nicht, was du willst,
was du wirklich denkst, was du fühlst. Ich weiss nicht, wie das mit uns weitergehen soll - .»

Ganz schön politisch oder mindestens philosophisch, hatte sich Louis
gedacht. Giorgia hatte ihn dabei seltsam eindringlich angesehen. Ihre wiederholten Vorträge, 
die, seiner Meinung nach, an den Haaren herbeigezogenen Vergleiche, hatten in diesem Moment 
das Fass zum Überlaufen gebracht. Vielleicht war es auch einfach ihr Blick gewesen, 
den Louis als nüchterne Distanz deutete, die ihn erschrocken den Kopf schütteln und sagen liess, er verstehe sie nicht. Selten so dankbar für seine Arbeit, hatte er sich rasch verabschiedet und ins Ristorante davongemacht. 

\qrblock{ILaria \emph{«Giudizi universali»}}{media/QR-Codes/040_Giudizi_universali_Ilaria_Spotify.png}

Louis setzte sich neben eine Schwarznase. Das Tier lag gemütlich
zwischen zwei mächtigen Steinen. Ein schräg darüber liegender
Gesteinsbrocken bot Sonnenschutz. Er steckte den leeren Sandwichbehälter zurück in den Rucksack,  legte sich auf den Rücken und schloss die Augen. Sein blosser Arm berührte das Fell des Schafes. Es kitzelte. 
Das vergangene Jahr mit Giorgia war kompliziert gewesen. Ihm war, als
sehe er ihre letzten, gemeinsamen Monate nur mehr in fahlem Licht. Hatte er sie womöglich zu einer Frau hochstilisiert, die sie gar nicht war? Und sich dabei mehr und mehr selbst verloren? Wie mochte es ihr gehen? Lag sie noch immer im Tiefschlaf, wurde sie  operiert? Wer wachte über sie?
Louis fühlte sich schwer und schwerer. Er rückte näher an das Schaf heran. Es öffnete für einen Moment die Augen und sah Louis direkt an. Der
Blick hatte etwas Tröstliches. Schliesslich fielen sie nebeneinander in einen kurzen Mittagsschlaf.
